\section{\texorpdfstring{$\Ham(M,\omega)$}{Ham(M, omega)} 的霍费尔拓扑}
\stfhIndex{I162}{huo^fei^er^tuo^pu}{霍费尔拓扑（Hofer topology）}
\label{sec:hofer-topology-ham}

霍费尔 (H90)、(H93) 引入的卓越哈密顿微分同胚伪范数定义为
\begin{equation}
 \|\phi\|=\inf_{H\mapsto\phi}\|H\|
 =\inf_{H\mapsto\phi}\operatorname{leng}(\phi_H).
 \label{eq:hofer-pseudonorm}
\end{equation}
下面命题的证明是一个涉及哈密顿代数的有趣练习。

\begin{proposition}\label{prop:hofer-pseudonorm-properties}
设 $\phi,\psi\in\Ham^c(M,\omega)$。则下列结论成立：
\begin{enumerate}[label=(\arabic*)]
 \item （对称性）$\|\phi\|=\|\phi^{-1}\|$；
 \item （三角不等式）$\|\phi\psi\|\leq\|\phi\|+\|\psi\|$；
 \item （辛不变性）对所有 $h\in\Symp(M,\omega)$，
 $\|h\phi h^{-1}\|=\|\phi\|$。
\end{enumerate}
\end{proposition}

\begin{proof}
回忆 (2.3.22) 中的逆哈密顿量公式
\[
 \overline H(t,x)=-H(t,\phi_H^t(x)),
\]
它生成 $(\phi_H^t)^{-1}$。显然有
\[
 \max_x\overline H_t=-\min_xH_t,
 \qquad
 \min_x\overline H_t=-\max_xH_t.
\]
因而 $\|\overline H\|=\|H\|$。此外，我们知道两次取“横线”运算是恒等运算。对所有 $H\mapsto\phi$ 在
$\|\overline H\|=\|H\|$ 两边取下确界，便由二者合并得到 (1)。

对于 (2)，回忆复合哈密顿量 $H\#F$ 定义为
\stfhIndex{I019}{00-symbol-019}{$H\#F$}
\[
 H\#F(t,x)=H(t,x)+F\bigl(t,(\phi_H^t)^{-1}(x)\bigr),
\]
当 $H\mapsto\phi$ 且 $F\mapsto\psi$ 时，它生成 $\phi\psi$。后一事实特别蕴含
\begin{equation}
 \|\phi\psi\|\leq
 \inf_{H\mapsto\phi,\,F\mapsto\psi}\|H\#F\|.
 \label{eq:hofer-composition-infimum}
\end{equation}
另一方面，有
\[
 \max_x(H\#F)_t\leq\max_xH_t+\max_xF_t,
 \qquad
 -\min_x(H\#F)_t\leq-\min_xH_t-\min_xF_t.
\]
于是得到
\begin{equation}
\begin{aligned}
 \|H\#F\|
 &=\int_0^1\left(\max_x(H\#F)_t-\min_x(H\#F)_t\right)dt\\
 &\leq\int_0^1\left(
   \max_xH_t+\max_xF_t-\min_xH_t-\min_xF_t\right)dt\\
 &=\|H\|+\|F\|.
\end{aligned}
\label{eq:hofer-composition-bound}
\end{equation}
合并\eqref{eq:hofer-composition-infimum}与\eqref{eq:hofer-composition-bound}，即证 (2)。

为证明 (3)，先回忆：对给定哈密顿量 $H$，拉回
$h^*H=H\circ h$ 对任意辛微分同胚 $h:M\to M$ 都生成
$h^{-1}\circ\phi_H^t\circ h$。显然有
$\|H\|=\|H\circ h\|=\|h^*H\|$。对所有 $H\mapsto\phi$ 取下确界，即证 (3)。
\end{proof}

这个范数随即在 $\Ham(M,\omega)$ 上给出一个自然的度量拓扑，称为
$\Ham(M,\omega)$ 的\emph{霍费尔拓扑}。这就是 (H90)、(H93) 中的原始定义。

现在沿用 (OhM07) 的论述，从
$\mathcal P^{\mathrm{ham}}(\Symp(M,\omega),\operatorname{id})$ 上的霍费尔拓扑出发，给出另一种更具概念性的方式来定义 $\Ham(M,\omega)$ 上的霍费尔拓扑。先重写 $\Ham(M,\omega)$ 的定义。以
$\mathcal P^{\mathrm{ham}}(\Symp(M,\omega),\operatorname{id})$ 表示始于单位元的哈密顿路径集合，并以
\begin{equation}
 \operatorname{ev}_1:
 \mathcal P^{\mathrm{ham}}(\Symp(M,\omega),\operatorname{id})
 \longrightarrow\Symp(M,\omega);
 \qquad \operatorname{ev}_1(\lambda)=\lambda(1)
 \label{eq:evaluation-map-ham-paths}
\end{equation}
表示评价映射。于是由 $\Ham(M,\omega)$ 的定义有
\stfhIndex{I020}{00-symbol-020}{$\Ham(M,\omega)$}
\begin{equation}
 \Ham(M,\omega)=\operatorname{ev}_1\bigl(
 \mathcal P^{\mathrm{ham}}(\Symp(M,\omega),\operatorname{id})\bigr).
 \label{eq:ham-as-evaluation-image}
\end{equation}
这个表达式以一种更明显的方式给出自然拓扑。

\begin{proposition}\label{prop:hofer-topology-quotient}
$\Ham(M,\omega)$ 上的霍费尔拓扑等价于使评价映射
\eqref{eq:evaluation-map-ham-paths}连续的最强拓扑。
\end{proposition}

\begin{proof}
由 $\Ham(M,\omega)$ 上霍费尔拓扑的定义，容易看出映射
$\operatorname{ev}_1^{\mathrm{ham}}$ 连续。因此，霍费尔拓扑弱于这一最强拓扑。反过来，只需证明 $\Ham(M,\omega)$ 在最强拓扑中的开集在霍费尔拓扑中也开。

设 $\mathcal U\subset\Ham(M,\omega)$ 在最强拓扑中开，并且
$\phi\in\mathcal U\subset\Ham(M,\omega)$。取 $H\mapsto\phi$。由最强拓扑的定义可知，
$\operatorname{ev}_1^{-1}(\mathcal U)$ 是
$\mathcal P^{\mathrm{ham}}(\Ham(M,\omega),\operatorname{id})$ 中
$\phi_H$ 的一个开邻域。因此，回忆
\stfhIndex{I012}{00-symbol-012}{$d(\phi',\phi)$}
\[
 d(\phi',\phi)
 =\inf_{H',H}\{d(\phi_{H'},\phi_H)\mid
   H'\mapsto\phi',\ H\mapsto\phi\},
\]
存在 $\epsilon>0$，使得只要
$d(\phi_{H'},\phi_H)<2\epsilon$，就有
$\phi_{H'}\in\operatorname{ev}_1^{-1}(\mathcal U)$，即
$\phi_{H'}^1\in\mathcal U$。现在断言，如果
$d_{\mathrm{Hofer}}(\phi',\phi)<\epsilon$，则 $\phi'\in\mathcal U$；这将表明 $\mathcal U$ 在霍费尔拓扑中开。依定义，可以找到 $F$，使得
$\phi_F^1=\phi'$ 且
\[
 d(\phi_F,\phi_H)<2\epsilon,
\]
这说明 $\phi_F\in\operatorname{ev}_1^{-1}(\mathcal U)$，特别地，
$\phi'\in\mathcal U$。证明完毕。
\end{proof}

现在的一个大问题是这一拓扑是否平凡。下一节将讨论这个问题。
